\section*{Chapter \ref{ch:s2:the-central-risk-book} --- The Central Risk Book}

\begin{solution}{exo:s2:the-central-risk-book:1}
\$2 million crosses; the net is a \$1 million purchase. Of the \$5 million gross, 80\% nets away.
\end{solution}

\begin{solution}{exo:s2:the-central-risk-book:2}
$2.326 \times 369\,000 \approx \$858\,000$; with the unrounded standard deviation, \$859\,000.
\end{solution}

\begin{solution}{exo:s2:the-central-risk-book:3}
$\$1\text{m} \times (0.0003 + 0.7 \times 0.02 \times \sqrt{1/50}) = \$1\text{m} \times (0.0003 + 0.00198) \approx \$2\,280$.
\end{solution}

\begin{solution}{exo:s2:the-central-risk-book:4}
Impact per dollar rises with the trade's size, so small trades are cheap per dollar. Netting removes the offsetting parts of the flows, and the net
that remains is traded as one larger order with more impact per dollar; the saving in cost is smaller than the share of volume crossed (19.5\%
against 47.0\%).
\end{solution}

\begin{solution}{exo:s2:the-central-risk-book:5}
Two reasons: later client flows offset part of what is kept, so it never needs trading; and what is traded goes in smaller daily pieces, with less
impact per dollar. The price is residual risk, as Garleanu and Pedersen's partial trading toward a target trades cost against risk.
\end{solution}

\begin{solution}{exo:s2:the-central-risk-book:6}
The common component adds up across desks and never nets. With more desks, the independent parts net to $\sqrt D$ times one desk's while the common
part grows as $D$, so it becomes a larger share of the net, and the net trades are larger with more impact per dollar.
\end{solution}

\begin{solution}{exo:s2:the-central-risk-book:7}
35\% a day: \$89\,300 a day, with a value at risk of \$449\,000. A tighter limit buys less risk at twice the cost of the \$1 million limit.
\end{solution}

\begin{solution}{exo:s2:the-central-risk-book:8}
At 5\% a day the carried P\&L has a standard deviation of \$684\,000 and a one-day 99\% value at risk of about \$1.59 million, far above the \$1
million limit. The cost falls to \$20\,300 a day, 93.9\% below desk by desk, because the book carries far more risk; the saving is a price for risk,
not a free improvement.
\end{solution}

\begin{solution}{pb:s2:the-central-risk-book:1}
\begin{enumerate}
\item A book that takes the risk desks acquire from clients, nets it and manages the net; pooling combines positions before hedging so offsets cancel.
\item Five desks, 100 stocks, 500 days, \$0.5 million flow sd per name a day with 10\% common variance; a three-factor model with 1.5\% idiosyncratic
volatility; 3 basis points half-spread and square-root impact with coefficient 0.7.
\item \$335\,000 a day desk by desk, \$270\,000 pooled (19.5\% less); 47.0\% of gross flow crosses.
\item Internalising a client's risk in anticipation of offsetting flow, or externalising it by hedging at once.
\item With independent flows the saving rises to 48.0\% with ten desks; with a common part it peaks at 19.7\% with six and falls to 17.1\% with ten.
\item It does not net, and grows faster than the netted part.
\item Temporary impact follows a 3/5 power of the trade rate rather than a square root; either way impact grows faster than size.
\item Impact per dollar is higher on the larger net trade than on the small offsetting pieces.
\item Hedging common factors at once with liquid instruments and leaving the rest to be traded out or offset.
\item Trade out a fixed share of inventory each day with futures on the beta-weighted rest: from \$269\,700 and no risk at 100\% to \$20\,300 and \$684\,000
of daily P\&L sd at 5\%.
\item It removes about a fifth of the carried risk for about \$500 a day.
\item 15\%: \$45\,300 a day and a value at risk of \$859\,000.
\item Pooling saves 19.5\% of the hedging cost; patient hedging within a \$1 million value-at-risk limit saves 86.5\%, carrying a daily P\&L standard
deviation of \$369\,000.
\item That each desk's hedge is executed as if alone in the market; in fact desks' orders meet in the same market.
\item Charge desks for the risk they hand over and credit them for the netting they provide, so they neither dump risk nor route around the book.
\item Offsetting flows stop coming, the inventory grows, and the rate must rise when trading is dearest.
\item Clients' information from one desk and trading by another; the book's limits and the desks' incentives.
\item Factor-hedged residual risk.
\item Chapter~24's desk internalises one client stream; the central book internalises across desks.
\item Netting and patience: the cost the bank does not pay for hedging risk that offsets or can wait.
\end{enumerate}
\end{solution}

\begin{solution}{iq:s2:the-central-risk-book:1}
A book that collects client risk from several desks, nets it and hedges only the remainder, factor risks first and the rest patiently. It saves the
spreads and impact of hedging offsetting risk and gives the bank one view of its exposures.
\end{solution}

\begin{solution}{iq:s2:the-central-risk-book:2}
Trace the cost-risk frontier on the book's actual flows, costs and volatilities, then pick the point that minimises cost plus a price of risk, or the
cheapest point within the risk limit; revisit when flows or volatility change.
\end{solution}

\begin{solution}{iq:s2:the-central-risk-book:3}
At a mid price at the time of the cross, with each desk charged the same, and a transfer charge reflecting the cost it would have paid externally
less a share of the saving; never at a price that favours one desk's clients over the other's.
\end{solution}

\begin{solution}{iq:s2:the-central-risk-book:4}
Value at risk and stress loss; factor exposures (market, sector, style); concentration by name against its liquidity; time to liquidate; limits on
inventory age; and escalation when offsetting flow dries up.
\end{solution}

\begin{solution}{iq:s2:the-central-risk-book:5}
Every desk's positions and trades, by instrument and time, streamed with consistent reference data and timestamps; netting by underlying; factor
exposures computed centrally; reconciliation with desks' books; and an audit trail of crosses.
\end{solution}

\begin{solution}{iq:s2:the-central-risk-book:6}
Desk by desk the expected linear cost is $D c\,E|q| = D c s\sqrt{2/\pi}$; pooled, the net has standard deviation $\sqrt D s$, so $c\sqrt D s\sqrt{2/\pi}$:
a factor $\sqrt D$ less. The impact cost is $k E|q|^{3/2} \propto s^{3/2}$ per desk, $D$ of them; pooled $\propto (\sqrt D s)^{3/2} = D^{3/4}s^{3/2}$:
a factor $D/D^{3/4} = D^{1/4}$ less.
\end{solution}
